\chapter{Tassonomia di Toda}
\label{ch:toda}
Il foglio ``Slay'' della tassonomia di Toda riporta 6 righe con solo 2
colonne compilate (elemento presente / esempio). Lo usiamo come metodo:
ogni riga dichiara cosa esiste oggi e cosa manca.

\begin{table}[h]
\centering\small
\begin{tabular}{@{}llp{6cm}@{}}
\toprule
Dimensione & Elemento Slay & Gap / colonne vuote \\
\midrule
Performance & Levels, Streaks & Manca curva level (level fisso 1), streak non tracciate \\
Social & Leaderboards (?) & Solo mock: nessuna classifica reale (cfr.\ Hub, Cap.~\ref{ch:hubteoria}) \\
Immersion & Avatar & Rogue Level 12 statico, nessuna personalizzazione \\
Economy & Hearts, Energy, Cards, Skills, XP, Coins & Valori fissi (vite 3/5, energia 5/5), nessuna spesa \\
Personalization & Custom path, Many choices & 1 sola roadmap navigabile su 3 \\
Feedback & Progress bar, XP notif, Correct answer feedback & Notifiche XP non collegate a HUD reale \\
\bottomrule
\end{tabular}
\caption{Le 6 righe Slay del foglio Toda con gap.}
\label{tab:toda}
\end{table}

\section{Metodo con gli esempi}
I fogli di esempio (Blank, Duolingo, EcoSteps) mostrano il procedimento:
(1) elencare gli elementi osservabili, (2) classificarli nelle 5 dimensioni
di Toda, (3) marcare esplicitamente assenze e dubbi (come il ``?'' sulle
leaderboard). Duolingo esemplifica streak/energy/campionati funzionanti;
EcoSteps mostra feedback ambientali minimali. Slay oggi sta tra i due:
ha il vocabolario (carte, cuori, energia) ma non ancora l'economia chiusa.
La Tab.~\ref{tab:toda} è quindi anche la checklist dei lavori futuri.
